\documentclass[../../../include/open-logic-section]{subfiles}

\begin{document}

\olfileid{sth}{card-arithmetic}{fix}
\olsection{نقاط $\aleph$ الثابتة}

قدمنا في~\olref[spine][]{chap} أن بديهية الاستبدال مفتقرة إلى مسوغ،
إذ تلزم عنها تراتبية شاهقة غاية الشموخ. وقد صار في وسعنا، بعد هذا القدر
من حساب الكاردينالات، أن نضرب لذلك الارتفاع مثلًا موجزًا.

فبيّن أن~$0 < \aleph_0$ و$1 < \aleph_1$ و$2 < \aleph_2$، وهلم
جرًا؛ بل إن تفاوت الحجم لا يزال يتعاظم في كل خطوة. وقد يدعو ذلك إلى
الظن بأن~$\kappa< \aleph_\kappa$ لكل عدد ترتيبي~$\kappa$.

غير أن هذا الظن باطل في~$\ZFC$. بل يمكننا إثبات وجود
\emph{نقاط ثابتة لـ$\aleph$}، أي كاردينالات~$\kappa$ تحقق
$\kappa=\aleph_\kappa$.

\begin{prop}\ollabel{alephfixed}
توجد نقطة ثابتة لـ$\aleph$.
\end{prop}

\begin{proof}
عرّف بالعود:
\begin{align*}
	\kappa_0 &= 0\\
	\kappa_{n+1} &= \aleph_{\kappa_n}\\
	\kappa&= \bigcup_{n < \omega}\kappa_n
\end{align*}
فيكون~$\kappa$ كاردينالًا بحسب
\olref[cardinals][classing]{unioncardinalscardinal}. ثم:
\[
	\kappa= \bigcup_{n < \omega} \kappa_{n+1} =
	\bigcup_{n < \omega}\aleph_{\kappa_n} =
	\bigcup_{\alpha < \kappa}\aleph_\alpha = \aleph_\kappa
\]
\end{proof}

وقد أفرد بولوس النقطة الثابتة لـ~$\aleph$ التي أنشأناها بمقالة. ولما
أثبت في صدرها وجود~$\kappa$ قال:
\begin{quote}
	[$\kappa$] عدد \emph{كبير جداً} بحسب معايير من لم يسبق لهم التعرض
	لنظرية المجموعات؛ وهو، فيما يبدو لي، كبير إلى حد يجعل صدق أي نظرية
	تتضمن بين أقوالها زعم وجود $\kappa$ من الأشياء على الأقل موضع شك.
	\citep[p.~257]{Boolos2000}
\end{quote}
ثم ختم مقالته أخيرًا بهذا السؤال:
\begin{quote}
	[هل] نشك في أنه، أياً يكن الحال في بداية القصة، حين نكون قد بلغنا هذا
	الموضع تكون العجلات قد أخذت تدور في فراغ، ولم نعد نصغي إلى وصف شيء
	واقع؟
	\citep[p.~268]{Boolos2000}
\end{quote}
فإن كنا قد جاوزنا حقًا «أي شيء واقع»، كان الاستبدال أولى ما يلام؛ فهو
البديهية التي بها تم البرهان. ويلزم بولوس حينئذ، فيما نظن، أن يعيد النظر
في قوله قبل ذلك بعقود إن الاستبدال لا يلزم عنه أثر «غير مرغوب فيه»؛
وانظر~\olref[replacement][extrinsic]{sec}.

ولكن أوجود~$\kappa$ من السوء بهذه المنزلة؟ يعين على وزن الأمر أن ننظر
في \emph{مفارقة تريسترام شاندي} عند راسل. كان تريسترام شاندي يدون أيام
حياته، ويستغرق تدوين يوم منها سنة كاملة. فهو كلما انقضت سنة ازداد
تأخرًا: بعد سنة لا يكون قد كتب إلا يومًا، وقد عاش 364 يومًا لم يكتبها؛
وبعد سنتين لا يكون قد كتب إلا يومين، وقد عاش 728 يومًا لم يكتبها؛ وبعد
ثلاث لا يكون قد كتب إلا ثلاثة أيام، وقد عاش 1092 يومًا لم يكتبها
\dots\footnote{نتجاهل السنوات الكبيسة.} ومع هذا، فإن كان تريسترام
خالدًا استطاع أن يكتب كل يوم، إذ يكتب اليوم~$n$ في السنة~$n$ من حياته؛
فتكتمل يومياته «عند نهاية الزمان».

فأي فرق بعيد بين هذا وبين أن يكون~$\alpha$ أصغر من
$\aleph_\alpha$---وأن يتفاقم الفارق بينهما---ثم نبلغ~$\kappa$ فيدرك
أحدهما الآخر ويصير~$\kappa = \aleph_\kappa$؟

ولندع هذا النظر، ولنسلم~$\ZFC$، ثم نختم بمزيد من خواص إنشاء النقاط
الثابتة. فالأحكام الثلاثة الآتية تبين، على جهة الحدس، وجود نقطة غير
تافهة تكون التراتبية عندها من العرض بمقدار ارتفاعها.

\begin{prop}\ollabel{bethfixed}
توجد نقطة ثابتة لـ$\beth$، أي يوجد $\kappa$ بحيث
$\kappa= \beth_\kappa$.
\end{prop}

\begin{proof}
كما في \olref{alephfixed}، باستعمال «$\beth$» بدلاً من «$\aleph$».
\end{proof}

\begin{prop}\ollabel{stagesize}
$\card{V_{\omega+\alpha}} = \beth_{\alpha}$. وإذا كان
$\omega \ordtimes \omega \leq \alpha$، فإن
$\card{V_\alpha} = \beth_\alpha$.
\end{prop}

\begin{proof}
أما الدعوى الأولى فباستقراء يسير ما فوق المتناهي. وتلزم الثانية من أن
$\omega \ordtimes \omega \leq \alpha$ يقتضي
$\omega + \alpha = \alpha$. ونستعين في إثباته بأحكام حساب الأعداد
الترتيبية من~\olref[ord-arithmetic][]{chap}. فأولًا
$\omega \ordtimes \omega = \omega \ordtimes (1 \ordplus \omega) =
(\omega \ordtimes 1) \ordplus (\omega\ordtimes\omega) = \omega
\ordplus (\omega \ordtimes \omega)$. فإذا كان
$\omega \ordtimes \omega \leq \alpha$، أي كان
$\alpha = (\omega\ordtimes\omega) \ordplus \beta$ لبعض $\beta$، فإن
$\omega \ordplus \alpha = \omega \ordplus
((\omega \ordtimes \omega) \ordplus \beta) =
(\omega \ordplus (\omega \ordtimes \omega)) \ordplus \beta =
(\omega \ordtimes \omega) \ordplus \beta = \alpha$.
\end{proof}

\begin{cor}
يوجد $\kappa$ بحيث $\card{V_\kappa} = \kappa$.
\end{cor}

\begin{proof}
خذ~$\kappa$ نقطة ثابتة لـ~$\beth$، وقد ضمن وجودها
\olref{bethfixed}. وواضح أن
$\omega \ordtimes \omega < \kappa$. فينتج من
\olref{stagesize} أن~$\card{V_\kappa} = \beth_\kappa= \kappa$.
\end{proof}

فالمراحل الواقعة دون~$V_\kappa$ بعدد ما في~$V_\kappa$ من
!!{element}. وعلى هذا يكون~$V_\kappa$، في التصور، عريضًا بقدر ارتفاعه.
ووجه الشبه بتريسترام شاندي قريب: ننتقل من مرحلة إلى التي تليها بأخذ
مجموعات القوى، فنعظم التراتبية في كل خطوة تعظيمًا بعيدًا؛ ثم يدرك عرضها
ارتفاعها «في النهاية»، أي عند المرحلة~$\kappa$.

وقد يقال: كم مرة يتساوى عرض التراتبية وارتفاعها؟ والجواب: بعدد الأعداد
الترتيبية. وبيان ذلك يحتاج إلى تفصيل.

عرّف حدًا~$\tau$ على الوجه الآتي. لأي~$A$، ضع:
\begin{align*}
	\tau_0(A) & \defis \cardsucc{\card{A}}\\
	\tau_{n+1}(A) & \defis \beth_{\tau_n(A)}\\
	\tau(A) & \defis \bigcup_{n < \omega}\tau_n(A)
\intertext{كما في \olref{bethfixed}، تكون $\tau(A)$ نقطة ثابتة لـ$\beth$
لكل $A$، ومن تعريف $\tau_0(A)$ يكون
$\card{A} < \tau(A)$ بوضوح. فانظر الآن إلى التعريف العودي الآتي:}
	W_0 &\defis 0\\
	W_{\alpha + 1} & \defis \tau(W_\alpha)\\
	W_\alpha & \defis \bigcup_{\beta < \alpha} W_\beta
	\text{، إذا كان $\alpha$ حدياً.}
\end{align*}
وهذا البناء متعين لجميع الأعداد الترتيبية. وكل~$W_\alpha$ كاردينال؛
وفي كل خطوة خلفية يكون
$W_\alpha=\card{W_\alpha}<\tau(W_\alpha)=W_{\alpha+1}$، وأما عند
المراحل الحدية فنأخذ الاتحاد. فتصير الخريطة
$\alpha\mapsto W_{\alpha+1}$ متزايدة قطعًا، فهي، بحسب التصور،
«!!a{injection}» من الأعداد الترتيبية إلى نقاط~$\beth$ الثابتة. وكما
تقدم، يكون~$V_{W_\alpha}$ عريضًا بقدر ارتفاعه، لكل~$\alpha$.

\end{document}
